\documentclass[border=10pt]{standalone}
\usepackage{fontspec}
\setmainfont{Noto Serif}
\newfontfamily\bengalifont{Noto Serif Bengali}[Script=Bengali]
\usepackage[Latin,Bengali]{ucharclasses}
\setTransitionsFor{Bengali}{\begingroup\bengalifont}{\endgroup}
\usepackage{newunicodechar}
\newunicodechar{।}{{\bengalifont\char"0964}}
\newunicodechar{॥}{{\bengalifont\char"0965}}

\usepackage[table]{xcolor}
\usepackage{booktabs,tabularx,array}
\definecolor{bapilrule}{HTML}{1F2937}
\definecolor{bapilstripe}{HTML}{F3F4F6}
\renewcommand{\arraystretch}{1.28}
\setlength{\tabcolsep}{5pt}

\begin{document}
\begin{minipage}{17cm}
{\large\bfseries সারণি ৪৮।}\ {\large কৌশলগত পর্যালোচনা-বিন্দু ও উদ্দীপক}\par\vspace{7pt}
\small
\rowcolors{2}{bapilstripe}{white}
\begin{tabularx}{\linewidth}{@{}>{\hsize=0.494\hsize\raggedright\arraybackslash}X>{\hsize=0.899\hsize\raggedright\arraybackslash}X>{\hsize=1.348\hsize\raggedright\arraybackslash}X>{\hsize=1.258\hsize\raggedright\arraybackslash}X@{}}
\arrayrulecolor{bapilrule}\toprule
\rowcolor{white}\textbf{কখন} & \textbf{পর্যালোচনা} & \textbf{মূল প্রশ্ন} & \textbf{সম্ভাব্য সিদ্ধান্ত} \\
\midrule
২০২৭-এর শেষে & প্রথম বছরের পর্যালোচনা (MS4) & উন্মুক্ত ডেটা দিয়ে P1 ও P2 কি মূল্য সৃষ্টি করেছে? শাসন কি কাজ করছে? & পাইলট অব্যাহত, পরিধি সমন্বয় বা বন্ধ; দ্বিতীয় ধাপের বাজেট নিশ্চিত \\
২০২৮-এর মাঝামাঝি & স্বাধীন মধ্যমেয়াদি মূল্যায়ন (MS6) & অর্থের যথার্থ ব্যবহার; গ্রহণ; আশ্বাসের মান & BPFM পর্যায় ১-এর সিদ্ধান্ত; দ্বিতীয় ধাপের জনবল \\
২০২৯-এর শেষে & ৩৬ মাসের পরিকল্পনার চূড়ান্ত মূল্যায়ন (MS9) & কোন সরঞ্জাম সংস্থাগুলোর মালিকানায় যাবে? প্রাতিষ্ঠানিক রূপ কি যথাযথ? & G4 হস্তান্তর; ব্লুপ্রিন্ট ২.০; বিকল্প গ ও ঘ-এর মধ্যে নির্বাচন \\
২০৩০ & পরবর্তী জাতীয় এআই নীতিচক্রের সঙ্গে সামঞ্জস্য & ২০৩০-পরবর্তী নীতিতে BAPIL কীভাবে সহায়তা করবে? & H3-এর জন্য সংশোধিত মিশন ও লক্ষ্য \\
২০৩৫ & দীর্ঘমেয়াদি কৌশলগত পর্যালোচনা & জাতীয় নেটওয়ার্ক কি যৌক্তিক? পদ্ধতিগুলো কি BAPIL-এর বাইরে ব্যবহৃত হচ্ছে? & নেটওয়ার্ক সম্প্রসারণ, একীকরণ বা পরিধি সংকোচন \\
২০৪১ & রূপকল্প পর্যালোচনা & BAPIL কি ২০৪১-এর পরিকল্পনার সঙ্গে সংগতিপূর্ণ জনমূল্য সৃষ্টি করেছে? & সনদ নবায়ন, সরকারে একীভবন বা সংরক্ষণাগার রক্ষা করে সমাপ্তি \\
যেকোনো সময় & ঘটনাভিত্তিক উদ্দীপক & নতুন আইন বা বিধি; গুরুতর ঘটনা; বড় প্রযুক্তিগত পরিবর্তন; অর্থায়ন হারানো & ৩০ দিনের মধ্যে স্টিয়ারিং বোর্ডের বিশেষ পর্যালোচনা \\
\bottomrule
\end{tabularx}
\par\vspace{3pt}{\scriptsize\color{gray} উৎস: BAPIL কারিগরি ব্লুপ্রিন্ট, সংস্করণ ১.০ (২০২৬)। CC BY 4.0।}
\end{minipage}
\end{document}
